\section{Method}
\label{sec:method}

\subsection{Streaming pipeline}
\Cref{fig:arch} shows the architecture. An ETL step converts each node's wide Parquet rows (one per 15 minutes, with avg/std/min/max of every metric and the Nagios label) into long events $(\text{node}, t, \text{metric}, \text{stat}, \text{value}, \text{state}, \text{anomaly flag})$. The metric list is configuration only; I use 12 metrics $\times$ 4 statistics = 48 values per node and window. A producer replays all nodes merged in event-time order at a configurable rate, serialises with Avro and a schema registry, and keys every message by node so that one partition preserves each node's order. The Kafka record timestamp is the publish time, which makes end-to-end latency measurable.

\begin{figure*}[t]
\centering
\begin{tikzpicture}[
  node distance=6mm,
  box/.style={draw, rounded corners=2pt, minimum height=10mm, minimum width=22mm, align=center, font=\footnotesize, fill=blue!15},
  topic/.style={box, fill=green!15},
  store/.style={box, fill=orange!18},
  engine/.style={box, fill=yellow!25},
  viz/.style={box, fill=purple!15},
  arr/.style={-{Stealth[length=2mm]}, thick}]
  \node[box] (data) {Data\\[-2pt]{\scriptsize M100 Parquet}};
  \node[box, right=of data] (etl) {ETL\\[-2pt]{\scriptsize reshape \& label}};
  \node[topic, right=of etl] (kafka) {Kafka\\[-2pt]{\scriptsize event stream}};
  \node[engine, above right=3mm and 9mm of kafka] (spark) {Spark\\[-2pt]{\scriptsize hourly aggregates}};
  \node[box, below right=3mm and 9mm of kafka] (det) {Online detectors\\[-2pt]{\scriptsize anomaly scores}};
  \node[store, right=21mm of $(spark)!0.5!(det)$] (db) {TimescaleDB\\[-2pt]{\scriptsize storage}};
  \node[viz, right=of db] (graf) {Grafana\\[-2pt]{\scriptsize dashboard}};
  \draw[arr] (data) -- (etl);
  \draw[arr] (etl) -- (kafka);
  \draw[arr] (kafka.east) -- ++(3.5mm,0) |- (spark.west);
  \draw[arr] (kafka.east) -- ++(3.5mm,0) |- (det.west);
  \draw[arr] (spark.east) -- ++(3.5mm,0) |- (db.west);
  \draw[arr] (det.east) -- ++(3.5mm,0) |- (db.west);
  \draw[arr] (db) -- (graf);
\end{tikzpicture}

\caption{Architecture. Two independent consumer groups read the same topic. Spark computes event-time aggregates and the detector service scores nodes; both write to TimescaleDB, which Grafana reads.}
\label{fig:arch}
\end{figure*}

\emph{Aggregation.} A Spark Structured Streaming job decodes the Confluent wire format (it strips the 5-byte header and fails if the schema id is not the expected one) and computes one-hour tumbling windows per (node, metric, stat) in event time, with a 30-minute watermark, i.e.\ two source samples. Results are written in update mode: every micro-batch is bulk-loaded into a staging table and merged with \texttt{INSERT \dots ON CONFLICT DO UPDATE}, so a restart that replays a batch produces the same rows.

\emph{Detection.} The detector service is a second consumer group on the same topic. It rebuilds the 48-value vector of each node and timestamp from the individual events, releasing a timestamp when all nodes are complete, when a node has moved three windows past it, or when a node has been silent for 30 s, and forward-fills values missing from late events. Detectors score each vector and then learn from it (prequential evaluation); scores go to TimescaleDB and flags to a second topic.

\subsection{Detectors}
All detectors implement one interface, \texttt{step(vectors, states)} $\to$ scores, and run unchanged in the Kafka service and in a batch mode that reads Parquet (used for the 2.5-year evaluation). Each node has its own model and its own online threshold.
\begin{itemize}
  \item \textbf{Half-space trees} \cite{tan2011hst} from CapyMOA and from River (10 trees, depth 8, window 96), on min-max scaled features, with a variant that feeds each node's deviation from the cluster median.
  \item \textbf{Relative entropy and Pearson correlation} after Rettig et al.\ \cite{rettig2015online}: the Kullback--Leibler divergence between value histograms of adjacent 16-window blocks, and the correlation of a node with the cluster median. I adopt the paper's thresholds (a multiple of the baseline mean, and mean plus $k\sigma$).
  \item \textbf{RUAD replication}: a per-node LSTM autoencoder (window of 10) on all 352 aggregates, trained without labels.
\end{itemize}
Thresholds of every detector are also calibrated on a validation period for the F1 comparison. I add three references: \emph{always OK}, the best \emph{single-feature rule}, and a supervised \emph{gradient-boosted trees} (GBT) model \cite{friedman2001gbm,pedregosa2011sklearn} trained on the current label as a ceiling.

\subsection{Drain anticipation}
Following GRAAFE \cite{molan2024graafe}, a node at time $t$ is positive when it is drained in $(t, t{+}h]$ with $h\in\{1\,\text{h},6\,\text{h}\}$; rows whose previous row is already positive are removed, so only drains \emph{entering} the horizon count (\cref{fig:task}). Features are the current 48 values, their changes over 1\,h and 6\,h, and deviations from the cluster and rack medians. The model is a histogram-based GBT trained on all positives and 5\% of negatives.

\begin{figure}[t]
\centering
\begin{tikzpicture}[x=1cm, y=1cm, font=\footnotesize]
  \draw[-{Stealth[length=2mm]}, thick] (0,0) -- (8.2,0) node[below left] {time};
  \fill[gray!35] (5,0) rectangle (7,0.5);
  \node at (6,0.25) {drained};
  \draw[decorate, decoration={brace, amplitude=3pt, mirror}] (3.6,-0.12) -- (4.95,-0.12)
    node[midway, below=3pt, align=center] {rows with\\$y=1$};
  \draw[decorate, decoration={brace, amplitude=3pt, mirror}] (5.05,-0.12) -- (7,-0.12)
    node[midway, below=3pt, align=center] {dropped\\(rising-edge rule)};
  \node[circle, fill=cmodel, inner sep=1.5pt, label=above:{features at $t$}] (t) at (0.9,0) {};
  \node[circle, fill=cmodel, inner sep=1.5pt, label=above:{$t{+}h$}] (th) at (2.3,0) {};
  \draw[{Stealth[length=1.6mm]}-{Stealth[length=1.6mm]}, cmodel] (0.9,0.75) -- (2.3,0.75)
    node[midway, above] {$h$};
  \node[align=left, anchor=west] at (0,1.55) {label known at $t{+}h$ $\Rightarrow$ the example is learned};
  \draw[dashed, gray] (3.6,0) -- (3.6,0.5) node[above, font=\scriptsize, gray] {$h$ before drain};
\end{tikzpicture}

\caption{Anticipation task and delayed labels. The label of the row at $t$ is only known $h$ windows later; the streaming model learns each example at that point.}
\label{fig:task}
\end{figure}

In the stream, features seen at $t$ are kept until $t{+}h$, when the live Nagios states reveal the label, and the example is then learned. I compare four policies for refitting the GBT on all resolved examples so far: never, every four weeks, and when a drift detector (ADWIN, DDM or OPTWIN) fires on the class-balanced prequential error, at most once a week. Incremental trees (River's adaptive random forest and Hoeffding adaptive tree, CapyMOA's adaptive random forest) serve as pure online alternatives.
